
\needspace{4\baselineskip}
\color{uniblue}{\subsection{Контроль отсутствия напряжения (КОН)}}
\color{black}

\begin{enumerate}

\item
КОН контролирует отсутствие напряжения на электрооборудовании. В устройстве контроль отсутствия напряжения представлен функцией «КОН».

\item
Контроль отсутствия напряжения на электрооборудовании срабатывает без выдержки времени при одновременном снижении всех линейных напряжений секции ниже уставки. Блокировка пусковых условий функции формируется при наличии неисправности цепей напряжения (БНН).

\item
Ввод функции в работу на этапе параметрирования устройства осуществляется программным переключателем SGF1 «Ввод\_функции». Информация о введенной функции отображается выходным сигналом «КОН: Ввод».

\item
Функция может быть выведена из работы оперативно путем активации сигнала «ОВ КОН», что характеризуется активным состоянием сигнала «КОН: Оперативный вывод» на выходе функции. 
\item
Алгоритм работы функции выполнен в соответствии с рисунком \ref{kon:img1}. В таблице \ref{kon:tbl1} приведены параметры, необходимые для настройки функции.

\vspace{3mm}
\begin{figure}[H]
\centering
\includegraphics[width=0.8\textwidth,height=0.8\textheight,keepaspectratio]{\fbpath/211. КОН ТН/_latex/img/kon.pdf}
\captionof{figure}{Функциональная схема логики КОН}\label{kon:img1}
\end{figure}

\end{enumerate}
